\section{Completion of the proof}\label{sec:completion}
\begin{proof}[Proof of Theorem~\ref{thm:main}]
The endpoint and symmetry reductions in Section~\ref{sec:preliminaries} reduce the problem to \(0<\rho<1\) and a nonconstant function. Theorem~\ref{thm:induction} proves Equation~\eqref{eq:main} when \(0<s\le .84\).
For larger \(s\), replace \(f\) by \(-f\) if necessary so that \(m\ge0\).
Proposition~\ref{prop:large-bias} gives the desired inequality whenever \(m\ge .74\).
If \(m<.74\) and the inequality were ever to fail, Lemma~\ref{lem:crossing} would provide a parameter satisfying Equation~\eqref{eq:crossing}.
At that parameter, either \(m\le r^2/2\) or \(m\ge r^2/2\).
Propositions~\ref{prop:central} and~\ref{prop:intermediate} exclude these possibilities unless \(f\) is a signed coordinate, for which equality is already known.
This proves the theorem for every \(n,f,\rho\).
\end{proof}


\subsection{The information-theoretic consequence}
For $0\le p\le1$, write
\[
 h_2(p)=-p\log_2p-(1-p)\log_2(1-p),
\]
with $0\log_2 0=0$. Mutual information in the next statement is measured
in bits. The argument is the implication of
Anantharam, Bogdanov, Chakrabarti, Jayram, and
Nair~\cite[Proposition~1 and Lemma~1]{ABCJN}, reproduced here in our
normalization. It applies to Boolean summaries of a uniform input,
including summaries with nonzero mean.

\begin{corollary}[Courtade--Kumar bound]\label{cor:ck}
Let $X$ be uniform on $\{-1,1\}^n$, let $f$ be Boolean, and let
$Y_i=X_iZ_i$, where the $Z_i$ are independent of $X$ and each other,
with $\Pr(Z_i=-1)=\varepsilon\in[0,1/2]$. Then
\[
 I(f(X);Y)\le1-h_2(\varepsilon).
\]
Signed coordinates attain equality.
\end{corollary}

\begin{proof}
Put $\rho=1-2\varepsilon$, $d=T_\rho f$, $m=\E f$,
$h=H(m)$, $b=\E H(d)$, and $s=H(\rho)$. Define
\[
 \varphi(t)=h_2\left(\frac{1-\sqrt{1-t^2}}2\right),\qquad 0\le t\le1.
\]
This function is increasing and convex. In fact, for $0<t<1$,
writing $v=\sqrt{1-t^2}$ gives
\[
 \varphi'(t)=\frac{t}{2v\log 2}\log\frac{1+v}{1-v}\ge0,
 \qquad
 \varphi''(t)=\frac{\log((1+v)/(1-v))-2v}{2v^3\log 2}\ge0.
\]
The last numerator is nonnegative because its derivative with respect
to $v$ is $2v^2/(1-v^2)\ge0$ and its value at zero is zero.
Continuity extends monotonicity and convexity to $[0,1]$.

The symmetric noise kernel gives
$\Pr(f(X)=1\mid Y)=(1+d(Y))/2$. The identity
$h_2((1+u)/2)=\varphi(H(u))$, valid for every $u\in[-1,1]$, and
Jensen's inequality therefore imply
\[
 I(f(X);Y)=\varphi(h)-\E\varphi(H(d))
 \le\varphi(h)-\varphi(b).
\]
Concavity of $H$ gives $b\le h$, and Theorem~\ref{thm:main} gives
$\delta:=h-b\le1-s$. For a convex function its increments over intervals
of fixed length are increasing with the left endpoint. Since $h\le1$,
\[
 \varphi(h)-\varphi(b)
 \le\varphi(1)-\varphi(1-\delta)
 \le\varphi(1)-\varphi(s)
 =1-h_2(\varepsilon).
\]
The same reasoning includes $\delta=0$ and the noise endpoints by
continuity. A signed coordinate is a uniform bit observed through a
binary symmetric channel, whose mutual information is
$1-h_2(\varepsilon)$.
\end{proof}

